\chapter{Feeling Restless}

\section*{Definition and Overview}
Feeling restless is a common and uncomfortable state of inner agitation in which a person finds it hard to sit still, relax, or concentrate. It may show as physical fidgeting, pacing, an inability to settle, or a racing mind that will not switch off. Restlessness is not a diagnosis in itself but a symptom that can arise from stress and anxiety, stimulant use, sleep deprivation, certain medicines, or medical conditions such as thyroid overactivity and restless legs syndrome. Most causes are treatable, and identifying the underlying driver is the key to relief. This chapter explains the common causes of restlessness, how it is assessed, and what helps.

\section*{Epidemiology}
Restlessness is very common and affects people of all ages. It is a core feature of anxiety disorders and attention deficit hyperactivity disorder (ADHD) in children and adults, and it is reported by a large proportion of people under significant stress. Restlessness is also a frequent side effect of medicines such as some antidepressants, asthma relievers, and decongestants, and of stimulants including caffeine, nicotine, and recreational drugs. Restless legs syndrome, a specific cause of nighttime restlessness, affects around 5--10\% of adults. Restlessness severe enough to disturb sleep or work is a common reason for medical consultation.

\section*{Aetiology and Pathophysiology}
\begin{itemize}
    \item \textbf{Stress and anxiety:} the body's fight-or-flight response, driven by adrenaline and cortisol, produces agitation, muscle tension, and an inability to settle
    \item \textbf{Stimulants:} caffeine, nicotine, and stimulant drugs overactivate the nervous system and cause jitteriness and restlessness
    \item \textbf{Medicines:} some antidepressants, corticosteroids, asthma relievers, and cold and allergy products can cause restlessness as a side effect
    \item \textbf{Sleep deprivation:} insufficient sleep leaves the nervous system over-aroused and undermines concentration and calm
    \item \textbf{Thyroid overactivity (hyperthyroidism):} excess thyroid hormone speeds up metabolism and the nervous system, causing restlessness, palpitations, and heat intolerance
    \item \textbf{Restless legs syndrome:} an irresistible urge to move the legs, mainly in the evening and at night, often with uncomfortable sensations
    \item \textbf{ADHD:} the brain's difficulty regulating attention and impulses produces chronic fidgeting and an inability to stay still
    \item \textbf{Withdrawal:} reducing or stopping alcohol, nicotine, or other substances can cause agitation and restlessness
    \item \textbf{Mental health conditions:} mania, agitated depression, and psychosis can all present with marked restlessness
    \item \textbf{Other medical causes:} low blood sugar, fever, pain, and some neurological disorders
\end{itemize}

\section*{Clinical Features}
\begin{itemize}
    \item An inner sense of agitation, tension, or being unable to relax
    \item Physical fidgeting, pacing, foot tapping, or an inability to sit still
    \item Difficulty concentrating, feeling easily distracted, and a racing mind
    \item Sleep disturbance, with difficulty falling or staying asleep
    \item Palpitations, sweating, and muscle tension, especially when driven by anxiety or stimulants
    \item An urge to move the legs at rest, relieved by movement, suggests restless legs syndrome
    \item Restlessness that appears suddenly, or is linked to a new medicine or substance, points to those causes
\end{itemize}

\section*{Differential Diagnosis}
The cause of restlessness is usually identified from the history, but several conditions overlap.
\begin{itemize}
    \item Generalised anxiety disorder and panic disorder
    \item ADHD in children and adults
    \item Hyperthyroidism
    \item Restless legs syndrome
    \item Medication and stimulant effects, including caffeine excess
    \item Alcohol or drug withdrawal
    \item Mood disorders, including bipolar disorder in a manic phase
    \item Sleep disorders, including insomnia and sleep apnoea
    \item Akathisia, a severe inner restlessness caused by certain antipsychotic medicines
\end{itemize}

\section*{When to Seek Medical Care}
Persistent restlessness that interferes with sleep, work, or daily life deserves assessment, as does restlessness that began with a new medicine. Restlessness in a child that affects school and behaviour should be discussed with a doctor.

\textbf{Call 000 (or local emergency services) for:}
\begin{itemize}
    \item Severe agitation with confusion, hallucinations, or behaviour that is dangerous to the person or others
    \item Chest pain, racing heart, or breathlessness with the restlessness
    \item Restlessness in someone who has taken an overdose of medicine or drugs
    \item Suicidal thoughts or behaviour
\end{itemize}

\section*{Diagnosis and Investigations}
\begin{itemize}
    \item \textbf{History:} onset, timing, triggers, medicines, caffeine and alcohol intake, sleep, stress, and mood
    \item \textbf{Physical examination:} pulse, blood pressure, tremor, and signs of thyroid overactivity
    \item \textbf{Blood tests:} thyroid function, blood sugar, and sometimes iron studies (for restless legs)
    \item \textbf{Mental health assessment:} screening for anxiety, depression, and ADHD
    \item \textbf{Sleep assessment:} history of leg urges at night, and sleep studies where sleep apnoea is suspected
    \item \textbf{Medication review:} checking all prescription, over-the-counter, and recreational substances
\end{itemize}

\section*{Management}
\begin{itemize}
    \item \textbf{Reduce stimulants:} cut caffeine, energy drinks, nicotine, and stimulant drugs, especially after midday
    \item \textbf{Address stress and anxiety:} relaxation techniques, breathing exercises, mindfulness, and psychological therapy
    \item \textbf{Improve sleep:} consistent bedtimes, a wind-down routine, and a dark, cool bedroom
    \item \textbf{Regular physical activity:} daily exercise uses up nervous energy and improves sleep, though vigorous exercise late at night can worsen restlessness
    \item \textbf{Review medicines:} with a doctor, adjust any medicine causing restlessness
    \item \textbf{Treat underlying conditions:} thyroid disease, ADHD, restless legs syndrome (with iron replacement where deficient, and medicines if needed), and anxiety disorders all have specific treatments
    \item \textbf{Limit alcohol:} alcohol fragments sleep and can worsen restlessness as it wears off
    \item \textbf{Structured routines:} for children and adults with ADHD, predictable routines and movement breaks reduce fidgeting
\end{itemize}

\section*{Complications}
\begin{itemize}
    \item Chronic insomnia and fatigue from an over-aroused nervous system
    \item Impaired concentration, work performance, and relationships
    \item Worsening of anxiety and mood disorders
    \item Increased risk of accidents and injury in severe agitation
    \item Missed diagnosis of conditions such as hyperthyroidism or restless legs syndrome
\end{itemize}

\section*{Prognosis and Outcomes}
The outlook is good once the cause is identified. Restlessness from stimulants resolves within days of reducing them; anxiety-related restlessness improves reliably with therapy and stress management; and thyroid disease, ADHD, and restless legs syndrome all respond well to specific treatment. Sleep and calm return as the underlying driver is addressed. For most people, a combination of cutting stimulants, improving sleep, and learning to manage stress restores a sense of physical and mental stillness within weeks.

\section*{Prevention and Self-Care}
\begin{itemize}
    \item Keep caffeine to the morning and limit energy drinks and other stimulants
    \item Build a regular sleep routine and aim for 7--9 hours of sleep
    \item Exercise regularly, and use movement breaks during long sedentary periods
    \item Practise slow breathing or mindfulness when agitation builds
    \item Check medicines and supplements for restlessness as a side effect
    \item Reduce alcohol and manage stress proactively
    \item See a doctor if restlessness persists, is severe, or is linked to other symptoms
\end{itemize}

\section*{Sources and Further Reading}
The following reputable sources provide further information for healthcare students and patients seeking reliable, current advice about feeling restless.
\begin{itemize}
    \item Healthdirect Australia. \textit{Anxiety and restless feelings.} https://www.healthdirect.gov.au/
    \item Beyond Blue. \textit{Anxiety support and information.} https://www.beyondblue.org.au/
    \item Restless Legs Syndrome Australia. \textit{RLS information.} https://www.rlsaustralia.org.au/
    \item NHS. \textit{Feeling restless and anxious.} https://www.nhs.uk/
    \item Mayo Clinic. \textit{Restless legs syndrome.} https://www.mayoclinic.org/
    \item MedlinePlus. \textit{Restlessness.} https://medlineplus.gov/
\end{itemize}
